\begin{table}[htbp]
\centering
\caption{Within-Group-65 Heterogeneity: Medications vs.\ Other Medical Supplies}
\label{tab:within_g65}
\begin{adjustbox}{max width=\textwidth}
\begin{threeparttable}
\small
\begin{tabular}{lcccc}
\toprule
 & \multicolumn{2}{c}{Log prices} & \multicolumn{2}{c}{Log firms} \\
\cmidrule(lr){2-3} \cmidrule(lr){4-5}
 & Base & PBU FE & Base & PBU FE \\
\midrule
\multicolumn{5}{l}{\textit{Panel A: Medications (class 6531)}} \\
$g65_{med} \times Pre$ & -0.1658*** & -0.1744*** & 0.1516*** & 0.1605*** \\
 & (0.0174) & (0.0172) & (0.0100) & (0.0101) \\
Observations & 573,080 & 573,080 & 678,547 & 678,547 \\
\midrule
\multicolumn{5}{l}{\textit{Panel B: Other medical supplies (non-6531)}} \\
$g65_{other} \times Pre$ & -0.0992*** & -0.0973*** & 0.0444*** & 0.0504*** \\
 & (0.0092) & (0.0086) & (0.0057) & (0.0057) \\
Observations & 593,920 & 593,920 & 705,616 & 705,616 \\
\bottomrule
\end{tabular}
\begin{tablenotes}
\small
\item \textit{Notes:} 18-month window. Panel A restricts group 65 to class 6531 
(medications, subject to CMED pharmaceutical price regulation), using all non-group-65 items 
as the comparison. Panel B restricts group 65 to non-medication medical supplies 
(hospital equipment, dental supplies, etc.). If the price effect were driven by 
differential pharmaceutical inflation rather than competition, it should appear 
primarily in Panel A. SE clustered at item level. 
* $p<0.10$, ** $p<0.05$, *** $p<0.01$.
\end{tablenotes}
\end{threeparttable}
\end{adjustbox}
\end{table}
